% Nouveau EGA V (cote 25): archive cote 25, site page 23 (PDF page 24).
% typescript. Draft compared: none.
% AI-transcribed by Claude Opus 5.5 (claude-opus-5-5, adaptive thinking, effort high), Message Batches API, batch ega25b (msgbatch_01A7n5YuXubUPNK1zsepi8RG), 2026-11-13; not yet checked by eye.
\scanpage{25/23}
\mspage{\hand{20}}

\uline{Proposition} 5.3.   Avec les notations précédentes, soit U l'ensemble des
$x\hand{\in} X_{\hand{\check{}}P}$ , d'image $\hand{\xi}$ dans $\hand{\check{}}P$ , tels que $\phi_{\hand{\xi}}$ soit $F_{k(\hand{\xi})}$-régulière en x. \var{\deleted{Suppo-%
sons de plus}}{}{deletion}{ega25_25_023-v0001} Alors:
% NOTE: the words "Si F de présentation finie et plat relativement à S," are encircled by hand, and a line leads from the loop down to a), between "a)" and "U".
a) \var{\inserted[typed]{\hand{Si} F de présentation finie et plat relativement à S\hand{,}}}{\hand{Si} F de présentation finie et plat relativement à S\hand{,}}{insertion}{ega25_25_023-v0002} U est ouvert, et $\underline{G}\hand{|}U$ est plat relativement à $\hand{\check{}}P$ .

b) Pour tout $s\hand{\in}\hand{\check{}}P$ , si $\hand{\eta}$ désigne le point générique de $\hand{\check{}}P_s$ , U con-%
tient $X_{k(\hand{\eta})}$.

Démonstration. a) Comme $F_{\hand{\check{}}P}$ est \var{\inserted[typed]{de présentation finie, et}}{de présentation finie, et}{insertion}{ega25_25_023-v0003} plat par rapport à $\hand{\check{}}P$ , la conclu-%
sion résulte de 11.3. (et même de $\underline{E}_{III}$ ... dans le cas S loc noeth).

b) On peut supposer S=Spec(k). Les cycles associés à $F_{k(\hand{\eta})}$ sont
\var{\inserted[typed]{\var{\xout{non vides}}{}{typewriter strike-out}{ega25_25_023-v0005}}}{\var{\xout{non vides}}{}{typewriter strike-out}{ega25_25_023-v0005}}{insertion}{ega25_25_023-v0004}
(en vertu du par.3 ) les images inverses des cycles \var{\inserted[typed]{$Z_i$}}{$Z_i$}{insertion}{ega25_25_023-v0006} associés à F. \var{\xout{En
vertu de 2.6, par exemple, si $Z_i$ est une partie sous-schéma fermé intègre
de X telle que}}{}{typewriter strike-out}{ega25_25_023-v0007} Si $f(Z_i)$ est fini, alors en vertu de 2.3. \var{\xout{X}}{}{typewriter strike-out}{ega25_25_023-v0008}
$Z_{i\,k(\hand{\eta})}\hand{\cap} Y = \emptyset$ , dans le cas contraire, on vertu de 2.6.\var{\inserted[typed]{par exemple}}{par exemple}{insertion}{ega25_25_023-v0009} on a
également $Z_{i\,k(\hand{\eta})}\hand{\cap} Y \hand{\neq} Z_{i\,k(\hand{\eta})}$ (raisons de dimensions; 2.3. déjà invo-%
qué dans 2.6. est sans doute une meulleure raison), ce qui prouve que $\phi_{\hand{\eta}}$
ne s'annule sur aucun des $Z_{i\,k(\hand{\eta})}$ , et prouve donc b).

\uline{Corollaire} 5.4. Soit V l'ensemble des $\hand{\xi}\hand{\in}\hand{\check{}}P$ tels que $\phi_{\hand{\xi}}$ soit \var{\xout{E}}{}{typewriter strike-out}{ega25_25_023-v0010} $F_{k(\hand{\xi})}$-régu-%
lier. \var{\inserted{Si F est de présentation finie}}{Si F est de présentation finie}{insertion}{ega25_25_023-v0011} \var{\over{A}{a}}{a}{writing over}{ega25_25_023-v0012}lors V est constructible et contient les points génériques des
fibres de $\hand{\check{}}P$ sur S, d'autre part si \var{\inserted{de plus}}{de plus}{insertion}{ega25_25_023-v0013} X est propre sur S\var{\inserted{et F plat sur S}}{ et F plat sur S}{insertion}{ega25_25_023-v0014}\hand{,} l'ensemble V est
ouvert. \var{\deleted{\hand{Si \unsure{$S=$ spec $(k)$|$S$ spec $(k)$} ($k$ un corps) alors en tous cas $f(\check{P}-V)$ est
de codimension $\unsure{\geq|\hand{\neq}} 2$.}}}{}{deletion}{ega25_25_023-v0015}

\var{\inserted[typed]{\uline{Remarque}}}{\uline{Remarque}}{insertion}{ega25_25_023-v0016}\var{\xout{Xxxxxxxxxx}}{}{typewriter strike-out}{ega25_25_023-v0017}5.5. \var{\xout{Supposons}}{}{typewriter strike-out}{ega25_25_023-v0018} soit $\hand{\xi}\hand{\in}\hand{\check{}}P$ au dessus de $s\hand{\in} S$, et supposons
que $F_s$ soit sans cycles associés immergés, alors on voit aussitôt que
$\hand{\xi}\hand{\in} V$ (notation de 5.4.) signifie aussi que\var{\xout{xxxxx}}{}{typewriter strike-out}{ega25_25_023-v0019} aucune composante
irréductible de supp $F_{k(\hand{\xi})}$ ne se trouve au dessus de $H_{\hand{\xi}}$ (et on peut
évidemment dans ce critère remplacer $k(\hand{\xi})$ par une extension arbitraire
de $k(\hand{\xi})$ ). Notons que l'hypothèse $(S_1)$ sur $F_s$ qu'on vient de faire es

\var{\deleted{\hand{$\times$ Ce dernier fait se vérifie par la méthode de 5.3. b) et la théorie
de la dimension (on remplace $\eta$ par $\xi$ tel que $\underline{O}_{\check{P},\xi}$ de dim $1$))}}}{}{deletion}{ega25_25_023-v0020}
\var{\margin{23}}{\placedmargin{23}}{marginal note placed}{ega25_25_023-v0021}
